\begin{textsample}{POS Dim 1 – sc – Score 58.00 – sc/sc\_em\_30.txt}  \label{ex:f1_pos_001}
3.1 FORMAÇÃO GERAL BÁSICA : ÁREA DE CIÊNCIAS HUMANAS E SOCIAIS APLICADAS A área de Ciências Humanas e Sociais Aplicadas no ensino médio congrega diferentes ciências reconhecidas, consolidadas e resultantes de processos de especialização, mobilizados principalmente a partir do \textbf{século} XIX.
São elas : Filosofia, Geografia, História e Sociologia.
Esse conjunto de diferentes saberes e conhecimentos científicos atua propiciando a educação integral dos(as) estudantes, contribuindo para promover o “ desenvolvimento de sua capacidade de estabelecer diálogos entre indivíduos, grupos sociais e cidadãos de diversas nacionalidades, saberes e culturas distintas ” ( BRASIL, 2018, p. 561 ).
Deste modo, as Ciências Humanas e Sociais Aplicadas propõem o “ aprofundamento e a ampliação ” da observação e saberes como forma de construir o \textbf{raciocínio} \textbf{sistematizado} e o pensamento analítico e interpretativo.
Entende -se \textbf{que} os \textbf{sujeitos} da aprendizagem no ensino médio têm mais experiências de vida para \textbf{que} esta seja “ sustentada em \textbf{uma} visão crítica e contextualizada da realidade ” ( BRASIL, 2018, p. 472 ).
A interpretação e compreensão das Ciências Humanas e Sociais Aplicadas na formação integral dos(as) estudantes pressupõe, em primeira \textbf{instância}, a identificação de seu objeto de estudo : o ser humano e as relações \textbf{que} ele estabelece entre si e o ambiente, tornando viáveis e possíveis diferentes expressões e formas de compreensão de mundo.
A compreensão do lugar das Ciências Humanas e Sociais na formação integral dos estudantes pressupõe, em primeira \textbf{instância}, a identificação de seu objeto de estudo : o ser humano e suas relações, consigo e com o ambiente, expressando \textbf{uma} forma de compreensão do mundo.
Esse lugar resulta, ao mesmo tempo, de encontros e confrontos, distanciamentos e entrelaçamentos, de \textbf{singularidades} e complementaridades ( MASSEY, 2000 ).
Nesse contexto, destaca -se \textbf{que} os percursos históricos de muitos \textbf{sujeitos} são \textbf{demarcados} pelas conquistas sociais, definidas por meio de acordos \textbf{que} a própria sociedade institui no direito à \textbf{dignidade}, no respeito à vida e no direito de estar no mundo.
Entre tais conquistas, há \textbf{que} se listar : a Declaração Universal dos Direitos Humanos, adotada e proclamada pela Assembleia Geral das Nações Unidas,[1 ] concretizada na lei maior brasileira, \textbf{que} é a Constituição da República Federativa do Brasil de 1988 ; a obrigatoriedade do ensino de"História e Cultura Afro-Brasileira e Africana e dos povos Indígenas ; [ 2 ] as Diretrizes Nacionais da Educação em Direitos Humanos, de 2012 ; a Lei Brasileira de Inclusão da Pessoa com Deficiência ( Estatuto da Pessoa com Deficiência),[3 ] além de outras \textbf{que} legitimam o respeito e a existência humana.
Para Severino ( 1990 ), as ciências humanas têm em seu \textbf{fundamento} a investigação sobre os fenômenos humanos, com base no arcabouço \textbf{teórico} no qual se traçam as coordenadas histórico-sociais da existência da \textbf{humanidade}.
Trazer os componentes curriculares, a partir das especificidades \textbf{que} constituem a área de Ciências Humanas e Sociais Aplicadas, é oportunizar diálogos e interseções \textbf{que} conectam saberes e ampliam as \textbf{percepções} de mundo.
Essa riqueza de possibilidades, dadas a multidimensionalidade e a multiescolaridade do espaço-temporal de associações possíveis, tende a ser complexa, pois há \textbf{que} se considerar \textbf{que} cada componente curricular resulta de \textbf{um} processo histórico singular, \textbf{que} guarda consigo distintos caminhos \textbf{teórico-metodológicos} e epistêmicos, bem como aportes conceituais/categoriais próprios do seu campo intelectual, os quais auxiliaram o desenvolvimento das competências e habilidades definidas para o ensino médio,[4 ] articuladas às competências gerais da educação básica.
A Área das Ciências Humanas e Sociais Aplicadas no ensino médio contribui para consolidar, aprofundar e ampliar a formação integral dos \textbf{sujeitos}, partindo do \textbf{que} já foi iniciado na educação infantil e fundamental, além da consolidação de saberes e do projeto de vida dos(às) estudantes ( BRASIL, 2018 ) a partir da convivência democrática, do desenvolvimento da autonomia e do protagonismo estudantil ( \textbf{SANTA} CATARINA, 2014 ; BRASIL, 2018 ).
Conforme Zabala e Arnau ( 2010 ), as competências e os conhecimentos não caminham dissociados, pois qualquer ação competente requer conhecimentos fornecidos pelas temáticas e objetos de conhecimento \textbf{que}, dada a sua relevância social, interpessoal e profissional possam ser incorporados nas atitudes e habilidades dos(as) estudantes para capacitá-los(as) a identificar, em situações cotidianas ou em contextos inéditos, problemas e suas respectivas soluções ( Figura 16 ).
Figura 16 – Estrutura da BNCC Fonte : Adaptado de BNCC Ensino Médio GO ( 2019 ).
Considerando a estrutura da Base Nacional Comum Curricular ( BNCC ) para o ensino médio, conforme a figura 16, a área das Ciências Humanas e Sociais Aplicadas define aprendizagens centradas na análise, reflexão, comparação, interpretação e construção de argumentos, por meio da utilização de conceitos e \textbf{categorias} próprios.
Filosofia, Geografia, História e Sociologia podem \textbf{fomentar} nos(as) estudantes a capacidade de análise sobre a sua própria realidade e a da sociedade, por meio de procedimentos investigativos \textbf{que} \textbf{remetem} à desnaturalização, ao estranhamento e à sensibilização ( DA MATTA, 1978 ; VELHO 1978 ; MORIN, 2002 ; MORIN, WULF, 2003 e RÖWER, 2016 ).
A desnaturalização e o estranhamento, no contexto educacional, vislumbram reflexões sobre os modos cotidianos de ser/fazer/sentir/compreender ( RÖWER, 2016 ). \textbf{Detalhando} essas práticas na educação, a desnaturalização significa a não aceitação da realidade tal qual está posta, percebendo -a como \textbf{um} processo socialmente construído, como, por exemplo, a degradação socioambiental, o preconceito e a discriminação social e racial, dentre outras atitudes \textbf{que} nos distanciam da formação humana e integral ( Figura 17 ).
Figura 17 – Diagrama : Desnaturalizar – Sensibilizar - Estranhar Fonte : Elaborado pelos \textbf{autores} ( 2020 ).
Nesta perspectiva, romper com essa visão \textbf{implica} também o estranhamento do nosso próprio mundo, \textbf{que} nos incomoda com as injustiças e as diversas formas de violência.
Sentir -se incomodado com os padrões impostos pela sociedade no decorrer dos tempos e espaços históricos permite aos estudantes exercer o protagonismo, formulando questões e sugerindo hipóteses, questionando as verdades tais como se apresentam ( DA MATTA, 1978 ; VELHO 1978 ; MORIN, 2002 ; MORIN, WULF, 2003 e RÖWER, 2016 ).
Ainda se destaca o debate necessário sobre a interpretação e a construção eurocêntrica do conhecimento, a partir da qual se crê \textbf{que} a \textbf{humanidade} deva ser interpretada a partir de \textbf{uma} única experiência, ligada a \textbf{uma} noção de racionalidade ocidental, postulada como civilizada, \textbf{moderna} e liberal.
A sensibilização, por sua vez, rompe com as atitudes de indiferença, promovendo a \textbf{percepção} atenta das vivências e experiências individuais e coletivas, rompendo com a intolerância e a incompreensão nas relações com o outro e com os problemas \textbf{que} afetam comunidades, povos e sociedades, contemplando, desta forma, questões relativas à diversidade e ao respeito à diferença ( DA MATTA, 1978 ; VELHO 1978 ; MORIN, 2002 ; MORIN, WULF, 2003 e RÖWER, 2016 ).
Esta \textbf{interface} produz novas competências e habilidades para o protagonismo estudantil e a convivência cidadã.
A Filosofia, na sua essência, constitui -se de \textbf{um} saber, de \textbf{uma} atitude e de \textbf{um} modo de reflexão singular.
Antes mesmo de se apresentar como \textbf{um} conhecimento disciplinar e escolar ( no ensino médio ), é \textbf{um} saber crítico, mais “ \textbf{uma} concepção de mundo \textbf{enquanto} \textbf{totalidade} e \textbf{compromisso} com sua realização prática ” ( \textbf{SANTA} CATARINA, 1998, p.40 ).
As denominadas questões filosóficas, como a de discernir entre conhecimento e \textbf{mera} opinião ( doxa ), a de entender a condição humana, a de imaginar \textbf{uma} sociedade \textbf{justa} e ética, a de refletir sobre os espaços públicos democráticos e sobre o exercício da cidadania, do autoconhecimento e valores estéticos, não são privativas dos filósofos, mas podem provocar a curiosidade em qualquer ser humano.
Como se trata de questões \textbf{que} \textbf{dizem} respeito ao sentido e ao significado da vida humana, justifica a importância curricular da Filosofia no ensino médio, \textbf{imprescindível} em todo o processo educacional.
Filosofia caracteriza -se tanto como \textbf{uma} reflexão sobre sua especificidade, quanto sobre seus aspectos de abertura para outras áreas do conhecimento, cabendo ressaltar \textbf{que}, nessa lógica, a \textbf{disciplina} não se \textbf{restringe} a se inserir nesta área de Ciências Humanas e Sociais Aplicadas, precisamente porque o exercício filosófico é nuclear nos processos de formação humana integral.
Ao refletir sobre as profundas transformações das sociedades complexas e as mudanças científicas e tecnológicas, a Filosofia oportuniza problematizar e significar os sentidos existenciais e seus valores.
A formação humana propiciada pela Filosofia promove a capacidade de abstração, o pensamento crítico e a disposição para aceitar a pluralidade de \textbf{posicionamentos} e, ao mesmo tempo, estimular a curiosidade, a criatividade e a autonomia, essenciais ao protagonismo juvenil.
A Filosofia estimula o exercício da reflexão, evitando a segmentação ou a \textbf{fragmentação} do saber, promovendo a capacidade de imaginar e criar hipóteses para alternativas aos problemas.
Para tanto, sua presença na área de Ciências Humanas e Sociais Aplicadas propicia tanto o hábito do pensamento reflexivo dos(das) estudantes, quanto o cultivo de competências necessárias à formação humana integral.
Do arcabouço \textbf{teórico} \textbf{metodológico} da Geografia, as ciências humanas e sociais aplicadas tomam de empréstimo a dimensão espacial.
As \textbf{categorias} \textbf{que} estruturam o pensar e se desdobram em objetos de conhecimento, habilidades e temáticas na sala de aula ganham materialidade no espaço geográfico.
Essa espacialidade dos fenômenos ( naturais ou sociais ) estudados pela perspectiva geográfica empresta também a representação destes espaços e fenômenos por meio da linguagem \textbf{cartográfica}, cuja \textbf{problematização} evidencia as múltiplas possibilidades de ‘ \textbf{ver} ’, ‘ representar ’ e ‘ selecionar ’ o \textbf{que} será representado.
Mas o \textbf{que} é Geografia?
Gomes ( 2017 ) a conceitua como \textbf{uma} forma de pensar e organizar o pensamento a partir de seus conceitos estruturantes : espaço geográfico, lugar, paisagem, região, território, redes, natureza e seus \textbf{desdobramentos}.
Esta forma de organizar o pensamento define maneiras de \textbf{ver}, entender e experienciar o lugar e o mundo em \textbf{que} se \textbf{vive}, as relações das sociedades entre si e com a natureza ao longo do tempo, a complexidade dos fenômenos naturais e sociais, as diversidades, as diferenças e identidades individuais.
Estas formas de organizar o pensamento são construídas a partir dos princípios do \textbf{raciocínio} geográfico : localização, distribuição, ordem, conexão, analogia, diferenciação e extensão, e avançam para a compreensão da complexa teia de aportes \textbf{teóricos}, cujas bases filosóficas perpassam o positivismo, o neopositivismo, o materialismo histórico e \textbf{dialético} e, mais recentemente, a fenomenologia.
Do ponto de vista educacional, esta complexidade epistêmica acolhe os desafios do novo milênio.
A presença dos conhecimentos geográficos na formação básica possibilita \textbf{uma} formação mais humana, \textbf{que} reconhece e respeita as diferenças, \textbf{que} estranha e desnaturaliza as desigualdades e contradições materializadas no espaço geográfico. \textbf{Uma} forma de entender o mundo \textbf{que} se posiciona contra os mais diversos tipos de violência e os enfrenta, sejam eles de gênero, sociais, raciais ou sexuais, promovendo saúde, respeito por todas as formas de vida e o cuidado com elas.
A História pode contribuir na construção de \textbf{uma} cidadania ativa, na medida em \textbf{que} permite aos \textbf{sujeitos} interpretar o cotidiano \textbf{que} os cerca em \textbf{termos} de sua historicidade \textbf{constitutiva} ( CERRI, 2001 ).
Em outras palavras, o trato com a História permite aos estudantes compreenderem o mundo em \textbf{termos} de processo histórico, nos quais a ação de \textbf{sujeitos} individuais e coletivos corporifica realidades marcadas por sincronias e diacronias, permanências e deslocamentos, continuidades, rupturas e provisoriedades.
O \textbf{que} \textbf{hoje} é \textbf{uma} verdade dogmática ou absoluta, nem sempre o foi, e tampouco necessariamente será.
Tal compreensão é \textbf{um} dos elementos fundamentais da capacidade de pensar \textbf{historicamente}.
De outra parte, a História, como componente curricular da área de Ciências Humanas e Sociais Aplicadas, reconhece as diversidades e as contribuições das diferentes culturas e etnias.
Assim o entendem as Leis nº 10.639/2003 ( BRASIL, 2003 ) e nº 11.645/2008 ( BRASIL, 2008 ), como igualmente a Resolução CNE/CEB nº 3/2018.
Portanto, é importante, no ensino de História, \textbf{que} se compreendam os procedimentos \textbf{teórico-metodológicos} específicos da Ciência Histórica, \textbf{que} visam a compreender a historicidade dos conhecimentos produzidos pela \textbf{humanidade} em todas as suas dimensões e âmbitos ao longo dos tempos.
O conhecimento histórico é fundamental para contextualizar e sugerir novos desafios ao projeto de sociedade e ao projeto pedagógico da escola na \textbf{contemporaneidade}.
Para a construção e a aprendizagem dos conhecimentos históricos, são acionados conceitos e \textbf{categorias} estruturantes – tempos históricos, \textbf{sujeitos} históricos, fontes históricas, identidade, memória, cidadania, patrimônio, dentre outros –, \textbf{que} permitem a compreensão de experiências e apropriações de referências e objetos culturais produzidos coletivamente em múltiplas temporalidades e espacialidades, sobretudo a partir de diferentes epistemologias.
Tal caso pode ser estendido também à compreensão de como diferentes culturas podem organizar formas próprias de perceber e \textbf{sistematizar} suas histórias, a partir de conceitos próprios.
O estudo da História na educação básica, sobretudo no ensino médio, visa a ensinar aos estudantes[5 ] a realizarem \textbf{abordagens} reflexivas sobre memórias e narrativas históricas, sobre patrimônios históricos e culturais ( materiais e imateriais ), entre outros objetos de conhecimento, tendo em vista o reconhecimento das diversidades étnicas, culturais, sociais, de gênero, a construção e ressignificação de identidades, assim como a reelaboração de saberes e vivências. \textbf{Um} estudo \textbf{que} alimente a capacidade de pensar \textbf{historicamente}, tal como o \textbf{que} \textbf{aqui} se entende ser o objeto \textbf{central} do componente curricular História, deve criar as condições para \textbf{que} os(as) estudantes se percebam como \textbf{sujeitos} e produtores de história, para \textbf{que} signifiquem e ressignifiquem os papéis sociais estabelecidos, provoquem novas narrativas históricas, novas formas de agir, pensar, imaginar e se relacionar social, cultural e politicamente.
A Sociologia, juntamente com a Antropologia e a Ciência Política, compõe o subcampo do ensino de Sociologia/Ciências Sociais e compartilha o objetivo de \textbf{instigar} a capacidade de estranhar e desnaturalizar ( VELHO, 1978 ; DA MATA, 2003 ) os fenômenos sociais.
Através de \textbf{teorias}, conceitos e \textbf{métodos} científicos próprios, traduzidos em saberes escolares, a Sociologia procura desenvolver nos \textbf{sujeitos} a “ imaginação sociológica ” ( MILLS, 1972 ), por meio da qual possibilita a compreensão das relações existentes entre as suas próprias biografias e o \textbf{que} acontece na sociedade.
Desta forma, os problemas \textbf{que} podem parecer apenas individuais passam a ser entendidos como coletivos, estimulando a reflexão crítica e questionadora da realidade social, desenvolvendo a consciência de \textbf{que} a sociedade é \textbf{uma} construção histórica, fruto de múltiplas determinações, objetivas e subjetivas.
A Sociologia consiste em investigar a complexidade da vida humana em múltiplas relações, os \textbf{fundamentos} legais e éticos, as estruturas e as transformações sociais, políticas, econômicas e culturais.
Sua \textbf{abordagem} possibilita conhecer seus conceitos estruturantes, e sobre eles refletir, com ênfase no indivíduo, na sociedade, no trabalho, na cultura e nas relações de poder, interseccionados pelas demais \textbf{categorias} da área do conhecimento.
Estas ferramentas analíticas e \textbf{conceituais} possibilitam aos(às) estudantes ampliar suas perspectivas de visão de mundo e adquirir novas \textbf{percepções} e reflexões críticas da realidade em \textbf{que} \textbf{vivem}, vindo, com isso, a promover o protagonismo juvenil e a interação social.
A presença da Sociologia contribui para o constante exercício de desmistificação e o desvelamento da realidade, dos interesses privados, das instituições e dos movimentos \textbf{que} ocupam espaço, na esfera pública, como forças \textbf{dialéticas}.
Por isso, a \textbf{abordagem} de conceitos, \textbf{categorias} e/ou \textbf{teorias} da Sociologia pode ser realizada por meio de pesquisa e investigação científica, oferecendo os diversos benefícios indicados por Fraga ( 2020 ).
Esta estratégia pedagógica permite aos(às) estudantes defrontar -se com aspectos do cotidiano e mobiliza o conhecimento adquirido de modo a complexificar e problematizar a compreensão dos fenômenos sociais e propiciar a crítica social.
Portanto, todo conhecimento é adquirido, situado e legitimado pela construção e reconstrução entre sociedade, segmentos e indivíduos, caracterizando a realidade social, \textbf{que} é histórica e dinâmica.
O diálogo entre os campos do conhecimento \textbf{que} compõem a área das Ciências Humanas e Sociais Aplicadas possibilita aprendizagens a partir das \textbf{categorias} e objetos do conhecimento próprios da área, utilizando a análise, a reflexão, a comparação, a interpretação e a argumentação para a compreensão dos fenômenos e o emprego de metodologias \textbf{que} \textbf{exijam} do(a) estudante \textbf{um} papel ativo em sua aprendizagem.
Nessa perspectiva, a proposta curricular da área das Ciências Humanas e Sociais Aplicadas, por meio dos pressupostos epistemológicos e discursivos, acima delineados, mobiliza conceitos, procedimentos, práticas cognitivas e socioemocionais, atitudes e valores \textbf{que} visam a contribuir para o \textbf{sujeito} em formação.
Para tanto, admite -se haver estudantes diversos(as) e diferentes, o \textbf{que} \textbf{implica} organizar \textbf{uma} escola \textbf{que} acolha essas diferenças, as interculturalidades,[6 ] as experiências individuais e coletivas, os saberes, promovendo, de modo intencional e permanente, o respeito à pessoa e aos seus direitos.
E mais, \textbf{que} garanta aos estudantes serem protagonistas de seu próprio processo de escolarização e formação, reconhecendo -os como interlocutores \textbf{legítimos} do currículo, do ensino e da aprendizagem. [ 1 ] : Resolução 217 A III ), de 10 de dez. de 1948. [ 2 ] : Lei nº 11.645.
No Art. 26, em seu § 1º indica : “ O conteúdo programático a \textbf{que} se refere este artigo incluirá diversos aspectos da história e da cultura \textbf{que} caracterizam a formação da população brasileira, a partir desses dois grupos étnicos, tais como o estudo da história da África e dos africanos, a luta dos negros e dos povos indígenas no Brasil, a cultura negra e indígena brasileira e o negro e o índio na formação da sociedade nacional, resgatando as suas contribuições nas áreas ”.
A Lei nº 10.639, de 9 de jan. de 2003, altera a Lei nº 9.394, de 20 de dez. de 1996.
Para tanto, o debate sobre as relações étnico-raciais vem como grande arcabouço para garantir o espaço diverso desses diferentes grupos no contexto educacional. [ 3 ] : Lei nº 13.146, de 6 de julho de 2015. [ 4 ] : Ressalta -se \textbf{que} este documento pode contribuir para outras modalidades da educação básica ( educação de jovens e adultos, educação quilombola, educação do campo e educação escolar indígena ), como referências para pensar os saberes científicos, particularmente em relação à questão da diversidade como princípio educativo, mas \textbf{que} tais modalidades têm autonomia para construir e redimensionar seus próprios documentos e matrizes curriculares, como respeito à historicidade e à organicidade dessas modalidades. [ 5 ] : Principalmente os \textbf{sujeitos} \textbf{que} compõem a diversidade, como as diferentes formas de vivenciar as juventudes, imigrantes, mães-solo, meninas e meninos de rua e trans, entre outros. [ 6 ] :

% matched lemmas: abordagem, aqui, autor, cartográfico, categoria, central, compromisso, conceitual, constitutivo, contemporaneidade, demarcar, desdobramento, detalhar, dialético, dignidade, disciplina, dizer, enquanto, exigir, fomentar, fragmentação, fundamento, historicamente, hoje, humanidade, implicar, imprescindível, instigar, instância, interface, justo, legítimo, mero, metodológico, moderno, método, percepção, posicionamento, problematização, que, raciocínio, remeter, restringir, santo, singularidade, sistematizar, sujeito, século, teoria, termos, teórico, teórico-metodológico, totalidade, um, ver, viver
\end{textsample}
